\subsection{Following the same algebraic branch to the collision}
The one-place construction distinguishes the chosen collision point
from the positive arithmetic branch. A continuous choice of that
same scalar branch gives a complementary obstruction.

\begin{finding}[\fid{F55}, Part 3: continuous genuine trace branches]
\label{f:continuoustracebranch}
Let $b>1$ and let $f:[1,\infty)\to\C$ be continuous.
Suppose one monic irreducible polynomial
$P(v,Y)\in\C(v)[Y]$, of positive $Y$-degree, satisfies
$P(v,f(v))=0$ wherever its coefficients are defined.
Suppose that, for every sufficiently large integer $n$,
\[
 f(b^n)=\sum_i m_i\alpha_i^n,
 \qquad m_i\in\Z_{>0},\quad \alpha_i\in\C^*.
\]
Then there are rational numbers $r_i$ such that
\begin{equation}\label{eq:continuoustracepositive}
 \alpha_i=b^{r_i},\qquad
 f(v)=\sum_i m_i v^{r_i}\quad(v\ge1),\qquad
 f(1)=\sum_i m_i>0.
\end{equation}
All real powers use their positive real branches.

Consequently the actual marked-$X$ trace law of
\fid{F54}, Part~3, is impossible for arbitrary mixed
base spectra if the chosen evaluations of the coefficient
family vary continuously along the positive parameter
axis and are regular at $v=1$ on that same choice.
No weight or full-centralizer condition is needed.
\end{finding}

\begin{proof}
By \fid{F54}, Part~1, there are a common $L>0$ and
integers $a_i$ with
\[
 \alpha_i^L=b^{a_i}.
\]
Put $p(t)=\sum_i m_i t^{a_i}$.
At the indices $n=Lk$ the trace law gives
$f(b^{Lk})=p(b^k)$.
The rational function $P(t^L,p(t))$ therefore has
infinitely many distinct zeros $t=b^k$, away from its
finitely many possible poles, and is identically zero.

The extension $\C(t)/\C(t^L)$ is cyclic Galois, with
automorphisms $t\mapsto\zeta t$, $\zeta^L=1$.
Since $P$ is irreducible and has $p(t)$ as a root,
its roots are among $p(\zeta t)$.
Indeed the orbit polynomial of $p(t)$ has coefficients
fixed by this Galois group, and its irreducible minimal
factor is $P$.
Thus, away from a finite exceptional set, for real $v\ge1$,
\begin{equation}\label{eq:finitepositivebranches}
 f(v)\in\{p(\zeta v^{1/L}):\zeta^L=1\}.
\end{equation}
The inclusion also holds at an exceptional point:
approach through nonexceptional values, pass to a
subsequence with one fixed $\zeta$ from this finite set,
and use continuity of both functions.

Set $p_0(v)=p(v^{1/L})$.
For every $\zeta$ and every positive $v$,
\[
 |p(\zeta v^{1/L})|
 \le\sum_i m_i v^{a_i/L}=p_0(v).
\]
If $p(\zeta v^{1/L})=p_0(v)$, equality in the triangle
inequality, with a positive real sum and strictly
positive coefficients, forces every summand to have
argument zero. Hence $\zeta^{a_i}=1$ for every $i$.
That branch is therefore identically $p_0$.
No genuinely different branch in
\eqref{eq:finitepositivebranches} meets $p_0$ anywhere
on the positive real axis. Negative or zero exponents
do not alter this argument.

The set
\[
 E=\{v\in[1,\infty):f(v)=p_0(v)\}
\]
is closed and contains every sufficiently large
$b^{Lk}$. It is also open: at a point of $E$, the
finitely many genuinely different branch values have
positive distance from $p_0$; continuity preserves
a smaller positive separation nearby, and
\eqref{eq:finitepositivebranches} forces $f$ to remain
on $p_0$.
Connectedness gives $E=[1,\infty)$.

Now $|\alpha_i|=b^{a_i/L}$, and for every sufficiently
large $n$,
\[
 \sum_i m_i\alpha_i^n
   =\sum_i m_i|\alpha_i|^n>0.
\]
Equality in the triangle inequality forces each
$\alpha_i^n$ to be positive real. Applying this at two
consecutive such indices forces each $\alpha_i$ itself
to be positive real. Taking $r_i=a_i/L$ proves
\eqref{eq:continuoustracepositive}.

For the actual module, $X_b$ is $\Delta(b)$ times the
projector onto the rank-eight $N(D-2)^4$ sector.
The assumed right-action commutation makes that sector
invariant under the invertible base operator $F_b$.
The continuous scalar branch
\[
 f(v)=\frac{\Tr_M(X(v)F(v))}{\Delta(b)}
\]
satisfies the irreducible minimal polynomial over
$\C(v)$ of this trace element in the fixed finite
coefficient field, normalized to be monic.
Its marked comparison values are the genuine power
traces of the eight-dimensional restriction of $F_b$.
The proved scalar theorem gives $f(1)=8$.
Regularity on the same branch and the actual matrix
identity $X(1)=0$ give $f(1)=0$, a contradiction.
In fact only commutation of $F_b$ with $X_b$ is needed
for this last trace calculation.
\end{proof}

The continuous choice is additional comparison data.
It is not supplied merely by a finite algebraic cover
or by choosing some place over $v=1$.
In the explicit fractional model, the arithmetic points
have $t=2^{n/2}>0$ over $v=t^2$, while the regular
collision is at $t=-1$.
The normalized marked function $g(t)=2(1+t)^2$ has
$g(1)=8$ and $g(-1)=0$.
Its marked-sector trace $h$ and operator entries have
poles at $t=1$, even though $g$ and the adjusted
ordinary total trace $56+g$ are regular there.
That model does not satisfy this same-branch regularity
condition. Its one-place assertion remains valid.

% Separate paragraph for Part III.
\paragraph{\fid{F55}, Part 3, in Grothendieck's language.}
The coefficient functions lie on a finite cover of the
parameter torus. The family, action and marked
endomorphism are affine coherent-module data in
EGA~I, Theorem~1.3.7, p.~86, Corollary~1.3.8, p.~87,
and Theorem~1.4.1, pp.~90--92.
The representation and commuting-morphism dictionary
is SGA~3, I, 4.7.2--3, with cote~90, scan~17 as the
audited action anchor \cite{EGA,SGA3,Archive}.
After $v=t^L$, the genuine trace is an effective sum
of torus characters $t^{a_i}$.
Its positive multiplicities are the effective cone in
$K_0$ of torus representations; the tensor-character
language is SGA~6, V, \S\S2--3 and cote~133, scan~18
\cite{SGA6,Archive}.
This is not a positivity assertion for the
$\gamma$-filtration or a conjectural canonical basis.
The continuous branch over the positive real locus
is extra comparison data on the complex points of
the cover; the EGA module/sheaf equivalence does not
provide it. The full separation and phase-removal
proof was supplied above, rather than attributed to
an archival theorem. A virtual alternating trace
does not acquire the required positive multiplicities
merely by having a specified numerical dimension.
